% Generated draft: requires mathematical and source review.
\documentclass[11pt]{article}
\usepackage[margin=25mm]{geometry}
\usepackage{fontspec}
\setmainfont{Latin Modern Roman}
\usepackage{amsmath,amssymb,mathtools,unicode-math}
\setmathfont{Latin Modern Math}
\usepackage{enumitem,needspace,xurl,hyperref}
\setlength{\parindent}{0pt}
\setlength{\parskip}{5pt}
\setlength{\emergencystretch}{3em}
\newcommand{\N}{\mathbb{N}}
\newcommand{\Z}{\mathbb{Z}}
\newcommand{\Q}{\mathbb{Q}}
\newcommand{\R}{\mathbb{R}}
\newcommand{\C}{\mathbb{C}}
\newcommand{\F}{\mathbb{F}}
\newcommand{\A}{\mathbb{A}}
\newcommand{\PP}{\mathbb{P}}
\newcommand{\E}{\mathbb{E}}
\newcommand{\eps}{\varepsilon}
\newcommand{\dd}{\,\mathrm{d}}
\begin{document}
\section*{UnsolvedMath ID: 800489}
\textbf{Source problem code:} catalogue entry
\subsection{Definitions and mathematical statement}
A closed symplectic manifold is a compact smooth manifold \(M\) without boundary equipped with a closed, nondegenerate differential two-form \(\omega\). A Hamiltonian diffeomorphism \(\varphi\colon M\to M\) is the time-one map of a smooth time-dependent Hamiltonian flow. A periodic point is a point \(x\in M\) such that \(\varphi^k(x)=x\) for some integer \(k\geq 1\); its least period is the smallest such integer.

\textbf{Counting convention requiring confirmation.} The supplied excerpt does not define \(P(T)\). One explicit candidate interpretation is to count distinct finite cycles of \(\varphi\), counting each cycle once, without counting its iterates as additional orbits. Under this interpretation, set
\[
\mathcal O(\varphi)=\bigl\{\{x,\varphi(x),\ldots,\varphi^{k-1}(x)\}:x\text{ has least period }k\bigr\},
\qquad
P_\varphi(T)=\#\{O\in\mathcal O(\varphi):|O|\leq T\},\quad T\geq 1.
\]
Here the count takes values in \(\{0,1,2,\ldots\}\cup\{+\infty\}\), with every infinite cardinality recorded as \(+\infty\), and \(\log(+\infty)=+\infty\). These conventions are an explicit proposed completion of the excerpt and are not supplied source definitions.

\textbf{Question under this candidate interpretation.} Is it true that for every closed symplectic manifold \((M,\omega)\) and every Hamiltonian diffeomorphism \(\varphi\colon M\to M\), either \(\mathcal O(\varphi)\) is finite, or
\[
\liminf_{T\to+\infty}\frac{\log P_\varphi(T)}{\log T}\geq 2?
\]
In the infinite-orbit alternative, \(P_\varphi(T)\geq 1\) for all sufficiently large \(T\), so the displayed expression is defined for sufficiently large \(T>1\) with the stated extended-real convention.

The supplied authors express a suspicion that the answer is positive, citing a lack of counterexamples, and announce two particular test cases that are absent from this record. This is an attributed expectation, not an established mathematical assertion. The title and provenance identify Problem 2, whereas the extracted statement is labelled Problem 1; the intended source item and counting convention require confirmation.
\subsection{Short English statement}
The question asks whether every Hamiltonian diffeomorphism of a closed symplectic manifold has either finitely many periodic orbits or periodic-orbit growth satisfying the displayed logarithmic lower bound of \(2\). The proposed interpretation counts distinct cycles by least period. The source excerpt leaves that convention unspecified and has conflicting problem numbers.
\subsection{Sources}
\begin{itemize}
\item Local UnsolvedMath registry: permanent ID \(800489\), source problem code \texttt{catalogue entry}, in the AMR Open Problem Lists collection. The supplied registry labels the record \texttt{open}; this status has not been independently verified.
\item Supplied source list: Albach, Barraud et al., \emph{Open Problems in Billiards and Quantitative Symplectic Geometry} (2026), \url{https://arxiv.org/abs/2602.12896}. The supplied provenance attributes the item to Bernhard Albach, Umberto Hryniewicz, and Alexandru Oancea and identifies it as Problem 2. The extracted statement instead begins with Problem 1.
\item The supplied provenance records access on 2026-10-16 and extraction by \texttt{arxiv\_latex}. This access information is supplied provenance, not a claim of independent source consultation. The supplied upstream revision is \texttt{d371ee917ac9647fbb06cd6fc529deec6f3dda82}.
\end{itemize}
\end{document}
